\section{TD1 --- Noyaux}

\begin{td-exo}[Couverture des arêtes par des cliques]\,\\
    On considère le problème suivant, appelé ECC.
    Étant donné un graphe \(G\) et un entier \(k\), le but est de déterminer s'il existe \(k\) cliques \(C_1,\ldots, C_k\) telles que pour chaque arête \(e\) du graphe, il existe \(i\) tel que \(e \in C_i\).

    Le but de l'exercice est de montrer l'existence d'un noyau pour ce problème avec au plus \(2^k\) sommets.
    Pour tout sommet \(u\in V(G)\), on dénote \(N(u) = \left\{ v, uv \in E(G) \right\}\) et \(N[u] = N(u) \cup \left\{ u \right\}\).
    On dit que \(u\) et \(v\) sont \defemph{jumeaux} si \(N[u] = N[v]\).
    On suppose dans tout l'exercice que le graphe n'a pas de sommets isolés.

    \begin{enumerate}
        \item Soit \(V\) un ensemble de taille \(2^k + 1\) et \(C_1,\ldots,C_k\) des sous-ensembles de \(V\).
        Montrer qu'il existe \(u, v \in V\) tels que pour tout \(i \in \left\{ 1,\ldots,k \right\}\), \(u \in C_i\) si et seulement si \(v \in C_i\).

        \item Montrer que si \(u\) et \(v\) sont jumeaux, alors les instances \((G, k)\) et \((G \setminus u, k)\) sont équivalentes pour le problème ECC.

        \item En déduire un algorithme de kernelisation avec au plus \(2^k\) sommets.
    \end{enumerate}

\end{td-exo}


\begin{td-sol}
    \begin{enumerate}
        \item On pose \(X_V \in {\left\{0, 1\right\}}^k\) le vecteur de taille \(k\) défini par \(X_V[i] = 1\) si \(v \in C_i\) et \(0\) sinon.
        Alors, on a \(2^k\) vecteurs différents de taille \(k\) et \(2^k + 1\) sommets.
        Par le principe des tiroirs, il existe donc \(u, v \in V\) tels que \(X_u = X_v\).

        \item Soit \(C_1,\ldots,C_k\) des cliques couvrant les arêtes de \(G\).
        On pose
        \begin{equation*}
            C'_i = \begin{cases}
                C_i \setminus \left\{ u \right\} & \text{si } u \in C_i,\\
                C_i & \text{sinon.}
            \end{cases}
        \end{equation*}
        Alors, \(C'_1,\ldots,C'_k\) sont des cliques couvrant les arêtes de \(G \setminus u\).
        Réciproquement, soit \(C'_1,\ldots,C'_k\) des cliques couvrant les arêtes de \(G \setminus u\).
        On pose
        \begin{equation*}
            C_i = \begin{cases}
                C'_i \cup \left\{ u \right\} & \text{si } v \in C'_i,\\
                C'_i & \text{sinon.}
            \end{cases}
        \end{equation*}
        Alors, \(C_1,\ldots,C_k\) sont des cliques couvrant les arêtes de \(G\).
    
        \item On peut supprimer les sommets jumeaux du graphe jusqu'à ce qu'il ne reste plus que \(2^k\) sommets.
        Ainsi, l'instance renvoyée par l'algorithme de kernelisation est équivalente à l'instance initiale et contient au plus \(2^k\) sommets.
    \end{enumerate}
\end{td-sol}

\begin{td-exo}[Décomposition en couronne : démonstration]\,\\
    Une \defemph{décomposition en couronne} d'un graphe \(G\) est une partition de ses sommets en trois ensembles \(C, H, R\) telle que :
    \begin{enumerate}[label=(\roman*)]
        \item \(C\) est non vide,
        \item \(C\) est un stable,
        \item il n'y a pas d'arêtes entre \(C\) et \(R\),
        \item l'ensemble des arêtes entre \(C\) et \(H\) contient un couplage de taille \(\left| H \right|\).
    \end{enumerate}

    \begin{enumerate}
        \item Montrer que si \(G\) a un sommet de degré 1 alors, il admet une décomposition en couronne.
    \end{enumerate}
    Soit \(G\) un graphe avec au moins \(3k + 1\) sommets, pour un certain \(k \in \bb N\) et sans sommets isolés.
    Le but de cet exercice est de montrer que \(G\) admet soit un couplage de taille \(k+1\), soit une décomposition en couronne.
    Soit \(M\) un couplage de \(G\) maximal pour l'inclusion \(\big(\)c'est-à-dire qu'il n'existe pas de couplage \(M'\) de \(G\) tel que \(M \subset M'\) et \(\left| M \right| < \left| M' \right|\big)\).
    On dénote \(V_M \subset V(G)\) l'ensemble des sommets incidents à une arête de \(M\).
    On dénote également \(I = V(G) \setminus V_M\).
    On construit un graphe biparti \(B = (V_M, I, E')\) où \(E' \subset E(G)\) est l'ensemble des arêtes incidentes à un sommet de \(V_M\) et un sommet de \(I\).

    % to have counter start at 2 use:
    % \setcounter{enumi}{1}
    % for local counter to start at 2 instead use:
    % 
    \begin{enumerate}[start=2]
        \item Montrer que \(I\) est un stable de \(G\).
    \end{enumerate}

    Soit \(X \subset V(G)\) un vertex cover de taille minimum dans \(B\).

    \begin{enumerate}[start=3]
        \item Montrer en utilisant le théorème de König que si \(X \cap V_M = \varnothing\) alors \(X = 1\) et il existe un couplage de taille au moins \(k + 1\) dans \(G\).

        \item On suppose maintenant que \(X \cap V_M \neq \varnothing\).
        On pose \(H = X \cap V_M, C = I \setminus X\) et \(R = V(G) \setminus (X \cup H)\).
        Montrer que \((C, H, R)\) est une décomposition en couronne de \(G\).
        (Indication : considérer \(M'\) un couplage de taille maximum dans \(H\) et appliquer le théorème de König.)

        \item Montrer qu'il existe un algorithme en temps polynomial, qui étant donné \(G\) calcule soit un couplage de taille au moins \(k + 1\), soit une décomposition en couronne.
    \end{enumerate}
\end{td-exo}

\begin{td-sol}
    Soit \(G\) un graphe et \(x\) un sommet tel que \(d_G(x) = 1\).
    On note \(y\) le voisin de \(x\).

    On définit les ensembles suivants :
    \begin{equation*}
        C = \{x\}, \quad H = \{y\}, \quad R = G \setminus \left\{x, y\right\}.
    \end{equation*}
\end{td-sol}

\begin{td-exo}[Décomposition en couronne : application au Vertex Cover]\,\\
    En utilisant la décomposition en couronne, construire un noyau pour Vertex Cover avec au plus \(3k\) sommets.\\
    \textit{Indication : généraliser la réduction vue en cours pour les sommets de degré 1.}
\end{td-exo}

\begin{td-exo}[Décomposition en couronne : application à SAT]\,\\
    Le but de l'exercice est de construire un noyau linéaire pour SAT basé sur la décomposition en couronne.
    Étant donné une formule sous forme normale conjonctive et un entier \(k\), le but est de décider s'il existe une affectation satisfaisant au moins \(k\) clauses.
    Soit \(\varphi\) une formule en forme normale conjonctive avec \(n\) variables et \(m\) clauses.
    \begin{enumerate}
        \item Montrer que si \(m\geq 2k\), alors l'instance est positive.
        
        \item Pour tout sous-ensemble \(X\) de variables, on note \(C(X)\) l'ensemble des clauses dans lesquelles au moins variable de \(X\) apparaît.
        Supposons qu'il existe un sous-ensemble non vide \(C\) de clauses tel que :
        \begin{enumerate}
            \item il existe un ensemble \(X\) de variables avec \(|X| > |C|\) tel que \(C(X) = C\) et
            \item pour tout \(X' \subsetneq X\), on a \(|C(X')| \geq |X'|\).
        \end{enumerate}
        Montrer en utilisant le théorème de Hall que pour toute affectation qui satisfait un nombre maximal de clauses, toutes les clauses de \(C\) sont satisfaites.

        \item Montrer grâce au théorème de Hall que tout graphe biparti admet une décomposition en couronne.

        \item Construire un noyau pour SAT avec au plus \(2k\) clauses et au plus \(k\) variables.
    \end{enumerate}
\end{td-exo}

\begin{td-exo}[Lemme du tournesol]\,\\
    On appelle \(d\)-HITTING le problème suivant.
    Soit \(\mathcal{A}\) une famille d'ensembles sur un univers \(U\) telle que chaque ensemble contient au plus \(d\) éléments et soit \(k\) un paramètre.
    Le but est de décider s'il existe un ensemble \(H \subset U\) de cardinal au plus \(k\) tel que pour tout \(S\in A\), on a \(S\cap H \neq \varnothing\).
    Noter que pour \(d=2\), le problème est équivalent à Vertex Cover.
    Le but de l'exercice est de construire un noyau pour \(d\)-HITTING en utilisation un lemme d'Erdős-Rado, aussi appelé le \emph{lemme du tournesol}.
    Un \defemph{tournesol à \(k\) pétales et de capitule \(Y\)} est une collection de \(k\) ensembles \(S_1,\ldots S_k\) telle que
    \begin{enumerate}[label=\roman*)]
        \item pour tout \(i\neq j\), on a \(S_i \cap S_j = Y\),
        \item les ensembles \(S_i \setminus Y\) (appelés pétales) sont non vides.
    \end{enumerate}

    \begin{lemma}[du tournesol]\,\\
        Soit \(\mathcal{A}\) une famille d'ensembles sur un univers U, telle que chaque ensemble de \(\mathcal A\) contient au plus \(d\) éléments.
        Si \(|\mathcal{A}| > d ! {(k-1)}^d\) alors \(\mathcal{A}\) contient un tournesol à \(k\) pétales.
    \end{lemma}\vspace{1mm}

    On pourra commencer par traiter la troisième question en admettant les résultats des deux premières.

    \begin{enumerate}
        \item Démontrer le lemme par récurrence sur \(d\).

        \item Montrer que le tournesol du Lemme du tournesol peut être calculé en temps polynomial en \(|\mathcal{A}|, |U|\) et \(k\).

        \item En déduire un noyau pour \(d\)-HITTING avec au plus \(d! k^d\) ensembles.
    \end{enumerate}
\end{td-exo}
